\documentclass[12pt,a4paper]{article}
\usepackage[utf8]{inputenc}
\usepackage{amsmath,amsfonts,amssymb}
\usepackage{listings}
\usepackage{xcolor}
\usepackage{geometry}
\usepackage{hyperref}
\usepackage{fancyhdr}

\geometry{margin=1in}
\pagestyle{fancy}
\fancyhf{}
\rhead{Python Lists - Complete Guide}
\lfoot{Python for Data Science}
\rfoot{\thepage}

% Python code styling
\lstset{
    language=Python,
    basicstyle=\ttfamily\small,
    keywordstyle=\color{blue}\bfseries,
    commentstyle=\color{green!60!black},
    stringstyle=\color{red},
    numbers=left,
    numberstyle=\tiny\color{gray},
    stepnumber=1,
    numbersep=8pt,
    showstringspaces=false,
    breaklines=true,
    frame=single,
    backgroundcolor=\color{gray!10}
}

\title{\textbf{Python Lists - Complete Guide}}
\author{Python for Data Science - Beginner's Level}
\date{\today}

\begin{document}

\maketitle
\tableofcontents
\newpage

\section{Overview}

Lists are one of the most versatile and commonly used data structures in Python. They are \textbf{ordered}, \textbf{mutable} (changeable), and allow \textbf{duplicate elements}. Lists can store elements of different data types.

\section{List Creation}

\subsection{Basic Creation}

\begin{lstlisting}
# Empty list
empty_list = []
empty_list = list()

# List with elements
numbers = [1, 2, 3, 4, 5]
mixed_list = [1, "hello", 3.14, True]
nested_list = [[1, 2], [3, 4], [5, 6]]

# Using list() constructor
from_string = list("hello")  # ['h', 'e', 'l', 'l', 'o']
from_range = list(range(5))  # [0, 1, 2, 3, 4]
\end{lstlisting}

\subsection{List Comprehension}

\begin{lstlisting}
# Basic comprehension
squares = [x**2 for x in range(10)]

# With condition
even_squares = [x**2 for x in range(10) if x % 2 == 0]

# Nested comprehension
matrix = [[i*j for j in range(3)] for i in range(3)]
\end{lstlisting}

\section{Accessing Elements}

\subsection{Indexing}

\begin{lstlisting}
fruits = ["apple", "banana", "cherry", "date"]

# Positive indexing (0-based)
print(fruits[0])    # "apple"
print(fruits[2])    # "cherry"

# Negative indexing (-1 is last element)
print(fruits[-1])   # "date"
print(fruits[-2])   # "cherry"
\end{lstlisting}

\subsection{Slicing}

\begin{lstlisting}
numbers = [0, 1, 2, 3, 4, 5, 6, 7, 8, 9]

# Basic slicing [start:end:step]
print(numbers[2:7])     # [2, 3, 4, 5, 6]
print(numbers[:5])      # [0, 1, 2, 3, 4]
print(numbers[5:])      # [5, 6, 7, 8, 9]
print(numbers[::2])     # [0, 2, 4, 6, 8]
print(numbers[::-1])    # [9, 8, 7, 6, 5, 4, 3, 2, 1, 0]
\end{lstlisting}

\section{Modifying Lists}

\subsection{Adding Elements}

\begin{lstlisting}
fruits = ["apple", "banana"]

# append() - Add single element at end
fruits.append("cherry")

# insert() - Add element at specific position
fruits.insert(1, "orange")

# extend() - Add multiple elements
fruits.extend(["grape", "mango"])

# Using + operator
more_fruits = fruits + ["kiwi", "peach"]
\end{lstlisting}

\subsection{Removing Elements}

\begin{lstlisting}
fruits = ["apple", "banana", "cherry", "banana", "date"]

# remove() - Remove first occurrence
fruits.remove("banana")

# pop() - Remove and return element by index
removed = fruits.pop(1)

# del statement
del fruits[0]

# clear() - Remove all elements
fruits.clear()
\end{lstlisting}

\section{List Methods}

\subsection{Search Methods}

\begin{lstlisting}
fruits = ["apple", "banana", "cherry", "banana", "date"]

# index() - Find first occurrence
idx = fruits.index("banana")

# count() - Count occurrences
count = fruits.count("banana")

# in operator - Check if element exists
print("apple" in fruits)     # True
\end{lstlisting}

\subsection{Sorting Methods}

\begin{lstlisting}
numbers = [3, 1, 4, 1, 5, 9, 2, 6]

# sort() - Sort in place
numbers.sort()

# sorted() - Return new sorted list
sorted_list = sorted(numbers)

# Custom sorting with key function
words = ["banana", "pie", "Washington", "book"]
words.sort(key=len)  # Sort by length
\end{lstlisting}

\section{Performance Considerations}

\begin{table}[h]
\centering
\begin{tabular}{|l|c|l|}
\hline
\textbf{Operation} & \textbf{Time Complexity} & \textbf{Notes} \\
\hline
Access by index & O(1) & Direct access \\
Search & O(n) & Linear search \\
Append & O(1) amortized & May need to resize \\
Insert at beginning & O(n) & Shifts all elements \\
Delete by index & O(n) & May shift elements \\
Sort & O(n log n) & Timsort algorithm \\
\hline
\end{tabular}
\caption{List Operations Time Complexity}
\end{table}

\section{Common Use Cases}

\subsection{Stack Implementation}

\begin{lstlisting}
stack = []

# Push
stack.append(1)
stack.append(2)
stack.append(3)

# Pop
top = stack.pop()  # Returns 3
\end{lstlisting}

\subsection{Matrix Operations}

\begin{lstlisting}
# 2D matrix using nested lists
matrix = [
    [1, 2, 3],
    [4, 5, 6],
    [7, 8, 9]
]

# Access element
print(matrix[1][2])  # 6

# Transpose matrix
transposed = [[row[i] for row in matrix] for i in range(len(matrix[0]))]
\end{lstlisting}

\section{Best Practices}

\begin{itemize}
    \item Use list comprehensions for simple transformations
    \item Avoid modifying lists while iterating over them
    \item Use appropriate data types for better performance
    \item Consider using \texttt{collections.deque} for frequent insertions/deletions at both ends
    \item Use \texttt{enumerate()} when you need both index and value
\end{itemize}

\section{Summary}

Lists are fundamental Python data structures that provide:
\begin{itemize}
    \item \textbf{Flexibility}: Store any type of data
    \item \textbf{Mutability}: Can be modified after creation
    \item \textbf{Order}: Maintain insertion order
    \item \textbf{Indexing}: Fast access by position
    \item \textbf{Rich methods}: Extensive built-in functionality
\end{itemize}

Master lists to build a strong foundation for Python programming and data science!

\end{document}